\documentclass[11pt,a4paper]{article}

\usepackage[margin=1in]{geometry}
\usepackage{amsmath,amssymb,amsthm}
\usepackage{booktabs}
\usepackage{array}
\usepackage{enumitem}
\usepackage[dvipsnames]{xcolor}

\theoremstyle{plain}
\newtheorem{lemma}{Lemma}
\newtheorem{proposition}[lemma]{Proposition}
\newtheorem{corollary}[lemma]{Corollary}
\theoremstyle{definition}
\newtheorem{fact}{Fact}
\newtheorem{statement}{Statement}

\newcommand{\PROVED}{\textcolor{ForestGreen}{\textbf{[PROVED]}}}
\newcommand{\TODO}{\textcolor{BrickRed}{\textbf{[TO PROVE]}}}
\newcommand{\NUM}{\textcolor{RoyalBlue}{\textbf{[VERIFIED NUMERICALLY]}}}
\newcommand{\Sig}{\mathcal{S}}
\newcommand{\Dis}{\mathcal{D}}
\newcommand{\Win}{\mathcal{T}}
\newcommand{\1}{\mathbf{1}}
\newcommand{\E}{\mathbb{E}}
\newcommand{\PP}{\mathbb{P}}

\title{\bf Gradient dynamics at zero initialisation\\[4pt]
\large Notation, proofs, and what remains}
\author{Working document --- ICLR / COLT paper}
\date{15 September 2026}

\begin{document}
\maketitle

\noindent
Every claim below carries one of three labels: \PROVED, \NUM\ (with the measurement quoted), or
\TODO. Nothing is asserted without one.

\tableofcontents
\bigskip\hrule

%=====================================================================
\section{The task}

The vocabulary is $\mathcal{V}=\Sig\cup\Dis$, of sizes $V_s$ and $V_d$:
\begin{itemize}[nosep]
\item \textbf{signal} symbols $\Sig=\{0,\dots,V_s-1\}$, which carry the Markov state;
\item \textbf{distractor} symbols $\Dis=\{V_s,\dots,V-1\}$, which carry nothing.
\end{itemize}

For \emph{each sequence independently} a transition table $\pi$ is drawn, assigning to every state
$s\in\Sig$ a successor set $\mathrm{succ}_\pi(s)\subset\Sig$ of size two.

\paragraph{How the successors are drawn.}
This matters for the proofs. From \texttt{data\_generation.py}, line 70:
\begin{verbatim}
order     = rng.random((B, v_sig, v_sig)).argsort(axis=2)
trans_arr = order[:, :, :2]   # first two of a random permutation of ALL of S
\end{verbatim}
so $\mathrm{succ}_\pi(v)$ is a uniform random $2$-subset of \emph{all} of $\Sig$, \emph{including
$v$ itself}.

\begin{fact}\label{fact:unif}
$\PP\bigl(w\in\mathrm{succ}_\pi(v)\bigr)=2/V_s$ for \emph{every} $w\in\Sig$, including $w=v$.
\end{fact}
\noindent\NUM\ $3000$ tables at $V_s=4$: the fraction of state-entries where the state is its own
successor is $49.95\%$, against $2/4=50\%$.

\begin{fact}\label{fact:sym}
The law of $\pi$ is invariant under any permutation (relabelling) of $\Sig$.
\end{fact}

Writing $\mathrm{st}(t)$ for the state after position $t$, the generative rule is
\begin{equation}
v_t=\begin{cases}
\text{uniform from }\mathrm{succ}_\pi\bigl(\mathrm{st}(t-1)\bigr),
& \text{w.p. }\alpha\quad\text{(state advances)},\\[3pt]
\text{uniform from }\Dis, & \text{w.p. }1-\alpha\quad\text{(state unchanged)}.
\end{cases}
\end{equation}
The target is $y_t=v_{t+1}$ and the loss is averaged over the scored window $\Win=[T/2,T)$.

\paragraph{Why the task is ``in-context''.}
Because $\pi$ is redrawn every sequence, no function of the weights alone beats the marginal. All
achievable performance must be read off the current context.

%=====================================================================
\section{Architecture and notation}\label{sec:arch}

There is no learned embedding; position $t$ is two concatenated one-hots,
\begin{equation}
x_t=\bigl[\,\underbrace{e_{v_t}}_{\text{symbol}}\;;\;\underbrace{e_t}_{\text{position}}\,\bigr]
\in\mathbb{R}^{d},\qquad d=V+T .
\end{equation}
Two attention layers use a \emph{strictly} causal mask ($u<t$, never $u=t$) and combined
query--key matrices $W^{(1)}\in\mathbb{R}^{d\times d}$, $W^{(2)}\in\mathbb{R}^{2d\times2d}$. The
value circuits are \emph{fixed} to the selector that passes the symbol one-hot.

\medskip\noindent\textbf{Layer 1.}
\begin{equation}
S^{(1)}_{t,u}=\tfrac{1}{\sqrt d}\,x_tW^{(1)}x_u^\top+M_{t,u},\qquad
a^{(1)}_{t,\cdot}=\mathrm{softmax}\bigl(S^{(1)}_{t,\cdot}\bigr),\qquad
h_t=\sum_{u<t}a^{(1)}_{t,u}\,e_{v_u}.
\end{equation}

\noindent\textbf{Layer 2}, reading $x^{(2)}_t=[\,x_t\,;\,h_t\,]\in\mathbb{R}^{2d}$:
\begin{equation}
S^{(2)}_{t,u}=\tfrac{1}{\sqrt d}\,x^{(2)}_tW^{(2)}x^{(2)\top}_u+M_{t,u},\qquad
a^{(2)}_{t,\cdot}=\mathrm{softmax}\bigl(S^{(2)}_{t,\cdot}\bigr),\qquad
f_t=\sum_{u<t}a^{(2)}_{t,u}\,e_{v_u}.
\end{equation}

\begin{quote}
\textbf{Structural point that drives everything.} $f_t$ is built from $e_{v_u}$, the key's
\emph{own symbol}, not from $h_u$. Hence $h$ influences the loss \emph{only through the scores}
$S^{(2)}$, never through the value path. This is exactly what makes Lemma~\ref{lem:step0}(i) work.
\end{quote}

\paragraph{Readout.}
A plain renormalisation against a fixed bias --- no unembedding, no logarithm, no softmax:
\begin{equation}\label{eq:readout}
p_t=\frac{f_t+b}{\sum_w\bigl(f_t[w]+b_w\bigr)},\qquad
b_w=\begin{cases}b_s,&w\in\Sig,\\ b_d,&w\in\Dis.\end{cases}
\end{equation}
Since $f_t$ is a convex combination of one-hots, $\sum_wf_t[w]=1$ \emph{identically}, so
\begin{equation}
\mathcal{Z}:=\sum_w\bigl(f_t[w]+b_w\bigr)=1+V_sb_s+V_db_d
\end{equation}
is a \emph{constant}, independent of the context. The value of $b_d$ is derived in
Section~\ref{sec:bias}.

%=====================================================================
\section{The state at zero initialisation}

Set $W^{(1)}=W^{(2)}=0$. Every score is $0$, so both softmaxes are uniform over $\{u<t\}$, giving
$a^{(1)}_{t,u}=a^{(2)}_{t,u}=1/t$ and therefore
\begin{equation}
h_t=f_t=m_t:=\frac1t\sum_{u<t}e_{v_u},
\end{equation}
the \emph{running symbol frequency} of the context. At zero initialisation the model \emph{is} the
in-context unigram predictor.

\medskip\noindent\NUM\ cross-entropy at initialisation $=2.039121$ at $\alpha=0.6$, matching the
closed form.

\begin{quote}
This is exactly the ``unigram phase'' that Edelman et al.\ (2024) observe mid-training. We do not
pass through it --- we start there.
\end{quote}

%=====================================================================
\section{Lemma 1: the gradient at initialisation \quad \PROVED}

\begin{lemma}[Step $0$]\label{lem:step0}
At $W^{(1)}=W^{(2)}=0$, with $\omega_t:=\bigl(m_t[y_t]+b_{y_t}\bigr)^{-1}>0$:
\begin{enumerate}[label=(\roman*),nosep]
\item \textup{(Gating.)} $\dfrac{\partial L}{\partial W^{(1)}}=0$ identically.
\item \textup{(Closed form.)}
\begin{equation}\label{eq:g2}
-\frac{\partial L}{\partial W^{(2)}}
=\frac{1}{|\Win|\sqrt d}\sum_{t\in\Win}\frac{\omega_t}{t}
\sum_{u<t}\Bigl(\1[v_u=y_t]-m_t[y_t]\Bigr)\,x^{(2)}_t\otimes x^{(2)}_u .
\end{equation}
\end{enumerate}
\end{lemma}

\subsection{Proof of (ii)}

\paragraph{Step 1: derivative with respect to $f$.}
With $L_t=-\log p_t[y_t]$ and $p_t[w]=(f_t[w]+b_w)/\mathcal{Z}$, treating $f$ as free,
\begin{equation}
g_t[w]:=\frac{\partial L_t}{\partial f_t[w]}=-\,\delta_{w,y_t}\,\omega_t+\frac{1}{\mathcal{Z}} .
\end{equation}
The first term comes from the numerator (only the target's coordinate appears there); the second
from the denominator (every coordinate appears once).

\paragraph{Step 2: push through the attention.}
Since $f_t=\sum_ua^{(2)}_{t,u}e_{v_u}$ we have $\partial L_t/\partial a^{(2)}_{t,u}=g_t[v_u]$.
The softmax Jacobian is $\partial a_{t,u'}/\partial S_{t,u}=a_{t,u}(\delta_{uu'}-a_{t,u'})$, so
\begin{equation}
\frac{\partial L_t}{\partial S^{(2)}_{t,u}}
=a^{(2)}_{t,u}\Bigl(g_t[v_u]-\sum_{u'}a^{(2)}_{t,u'}g_t[v_{u'}]\Bigr)
=a^{(2)}_{t,u}\bigl(g_t[v_u]-g_t\cdot f_t\bigr).
\end{equation}

\paragraph{Step 3: the bracket collapses.}
Using $\sum_wf_t[w]=1$,
\begin{equation}
g_t\cdot f_t=\sum_wf_t[w]\,g_t[w]=-\,m_t[y_t]\,\omega_t+\frac{1}{\mathcal{Z}},
\end{equation}
so the two $1/\mathcal{Z}$ terms cancel and
\begin{equation}\label{eq:bracket}
g_t[v_u]-g_t\cdot f_t=-\,\omega_t\bigl(\1[v_u=y_t]-m_t[y_t]\bigr).
\end{equation}

\paragraph{Step 4: chain to the weights.}
With $\partial S^{(2)}_{t,u}/\partial W^{(2)}=x^{(2)}_t\otimes x^{(2)}_u/\sqrt d$ and
$a^{(2)}_{t,u}=1/t$ at initialisation, assembling \eqref{eq:bracket} and flipping the sign for the
descent direction yields \eqref{eq:g2}. \qed

\begin{quote}
\textbf{Plain reading.} The coefficient on key $u$ is positive exactly when that key's symbol
\emph{is} the target. The first step says: \emph{attend more to positions showing the answer}.
\end{quote}

\subsection{Proof of (i)}

By Section~\ref{sec:arch}, $h$ enters the loss only through $S^{(2)}$. Differentiating,
\begin{equation}
\frac{\partial S^{(2)}_{t,u}}{\partial h_u}
=\tfrac{1}{\sqrt d}\bigl(x^{(2)}_tW^{(2)}\bigr)_{d:2d},
\qquad
\frac{\partial S^{(2)}_{t,u}}{\partial h_t}
=\tfrac{1}{\sqrt d}\bigl(W^{(2)}x^{(2)\top}_u\bigr)_{d:2d},
\end{equation}
and both carry $W^{(2)}$ as a \emph{multiplicative factor}, hence both vanish at $W^{(2)}=0$.
Every path from $W^{(1)}$ to $L$ runs through $h$, so $\partial L/\partial W^{(1)}=0$. \qed

\medskip\noindent\NUM\ $\lVert\partial L/\partial W^{(1)}\rVert=0.000\mathrm{e}{+}00$ exactly ---
zero, not small.

\begin{quote}
\textbf{Consequence.} The two layers \emph{cannot form simultaneously}. Layer 2 must move first.
\end{quote}

%=====================================================================
\section{The block decomposition}

Split $W^{(2)}$ into four $d\times d$ blocks matching $x^{(2)}=[x;h]$:
\begin{equation}
W^{(2)}=\begin{bmatrix}P&Q\\R&C\end{bmatrix},\qquad
\sqrt d\,S^{(2)}_{t,u}
=\underbrace{x_tPx_u^\top}_{\text{symbol}\times\text{symbol}}
+\underbrace{x_tQh_u^\top}_{\textbf{induction}}
+\underbrace{h_tRx_u^\top}_{\text{cross}}
+\underbrace{h_tCh_u^\top}_{\text{state}\times\text{state}} .
\end{equation}
Restricting to the signal sub-blocks and reading off \eqref{eq:g2}, for $v,w\in\Sig$:
\begin{align}
G_P[v,w]&\propto\sum_t\1[v_t{=}v]\,\frac{\omega_t}{t}
\sum_{u<t}\bigl(\1[v_u{=}y_t]-m_t[y_t]\bigr)\,\1[v_u{=}w],
\label{eq:GP}\\
G_Q[v,w]&\propto\sum_t\1[v_t{=}v]\,\frac{\omega_t}{t}
\sum_{u<t}\bigl(\1[v_u{=}y_t]-m_t[y_t]\bigr)\,m_u[w].
\label{eq:GQ}
\end{align}
The outer product $x^{(2)}_t\otimes x^{(2)}_u$ supplies $\1[v_t{=}v]$ on the query side; on the key
side it supplies $\1[v_u{=}w]$ for the $P$ block and $h_u[w]=m_u[w]$ for the $Q$ block.

\begin{center}
\begin{tabular}{@{}l>{\raggedright\arraybackslash}p{6.1cm}>{\raggedright\arraybackslash}p{5.4cm}@{}}
\toprule
block & rule a positive diagonal would encode & what must hold for it to help \\
\midrule
$G_P$ & ``when I am $v$, prefer keys whose \emph{symbol} is $w$''
      & a \emph{fixed} bigram table, stored in weights \\
$G_Q$ & ``when I am $v$, prefer keys whose \emph{preceding context} looks like $w$''
      & nothing fixed --- read from the context \\
\bottomrule
\end{tabular}
\end{center}

%=====================================================================
\section{Statement (a): the bigram block \quad \TODO}\label{sec:a}

\begin{statement}\label{st:a}
$\E\bigl[G_P[v,v]\bigr]-\E\bigl[G_P[v,w]\bigr]\le 0$ for all $w\ne v$, in expectation over the
table $\pi$ and the sequence.
\end{statement}

\subsection{What symmetry already gives}

By Fact~\ref{fact:sym}, a permutation of $\Sig$ fixing $v$ and swapping $w\leftrightarrow w'$ maps
the ensemble to itself, so
\begin{equation}
\E\bigl[G_P[v,w]\bigr]=\E\bigl[G_P[v,w']\bigr]\qquad\forall\,w,w'\ne v .
\end{equation}
All off-diagonal entries in a row are equal. \emph{This says nothing about the diagonal}: no
permutation fixes $v$ while moving $v$, so $\E[G_P[v,v]]$ is a genuinely separate quantity, and
nothing forces it to equal the others.

\subsection{The two terms}

With $N_t(w):=\bigl|\{u<t:v_u=w\}\bigr|$ the prefix count of $w$, \eqref{eq:GP} reads
\begin{equation}
G_P[v,w]=\sum_t\1[v_t{=}v]\frac{\omega_t}{t}
\Bigl[\;\underbrace{\1[w{=}y_t]\,N_t(w)}_{\text{term A}}
\;-\;\underbrace{m_t[y_t]\,N_t(w)}_{\text{term B}}\;\Bigr].
\end{equation}

\begin{description}
\item[Term A.] By Fact~\ref{fact:unif},
$\E\bigl[\1[w{=}y_t]\mid\mathrm{st}(t)=v\bigr]=\alpha/V_s$ for \emph{every} $w$, including $w=v$.
So term A contributes \emph{equally} to diagonal and off-diagonal.
\item[Term B.] This carries the asymmetry. Conditioned on $v_t=v$, the prefix count $N_t(v)$ is
\emph{elevated} relative to $N_t(w)$ --- the chain has just visited $v$, and states recur. The
baseline subtraction therefore bites \emph{harder} at $w=v$.
\end{description}

\begin{quote}
So the expected conclusion is not ``the diagonal is zero'' but ``the diagonal is strictly
negative'': the in-weights block is pushed \emph{away} from the identity.
\end{quote}

\medskip\noindent\NUM\ at $\alpha=0.6$, diagonal minus off-diagonal of the signal $4\times4$:

\begin{center}
\begin{tabular}{@{}lrl@{}}
\toprule
block & value & sign \\
\midrule
$G_P$ (bigram) & $-1.475\times10^{-4}$ & negative \\
$G_Q$ (induction) & $+9.621\times10^{-5}$ & positive \\
$G_C$ (state$\times$state) & $+9.804\times10^{-6}$ & positive \\
\bottomrule
\end{tabular}
\end{center}

\subsection{What the proof needs}

\begin{enumerate}[nosep]
\item Show term A cancels exactly between diagonal and off-diagonal. Immediate from
Fact~\ref{fact:unif}.
\item Show $\E\bigl[m_t[y_t]N_t(v)\bigr]-\E\bigl[m_t[y_t]N_t(w)\bigr]>0$: the prefix count of the
\emph{current} state exceeds that of any other symbol in expectation. This is a statement about
the stationary behaviour of the chain and should \emph{not} require a mixing bound --- it reduces
to comparing $\PP(v_s{=}v\mid v_t{=}v)$ against $\PP(v_s{=}w\mid v_t{=}v)$, summed over $s<t$.
\end{enumerate}

\begin{quote}
\textbf{Start here.} Self-contained, no mixing. A few hours will show whether the algebra opens up.
\end{quote}

%=====================================================================
\section{Statement (b): the induction block \quad \TODO\ (the hard one)}\label{sec:b}

\begin{statement}\label{st:b}
$\E\bigl[G_Q[v,v]\bigr]-\E\bigl[G_Q[v,w]\bigr]>0$ for all $w\ne v$.
\end{statement}

\subsection{The mechanism, concretely}

$G_Q[v,w]$ weights each key $u$ by $m_u[w]$ --- the frequency of $w$ in \emph{the key's own
prefix} --- and rewards keys whose symbol is the target. With $V_s=4$:

\begin{center}
\begin{tabular}{@{}lccccccc@{}}
\toprule
position & 0 & 1 & 2 & 3 & 4 & 5 & 6 \\
symbol   & 0 & 5 & \textbf{2} & 6 & 0 & 4 & \textbf{2} \\
state    & 0 & 0 & 2 & 2 & 0 & 0 & 2 \\
\bottomrule
\end{tabular}
\end{center}

Suppose at time $t$ the state is $v=0$ and the target is $y_t=2$. The gradient rewards keys with
symbol $2$: positions $2$ and $6$. \emph{Both were produced by a transition out of $0$.} That is
not coincidence, it is how the data is generated. So $m_u$ at those keys carries extra mass on
$0$, which is $v$ itself.

\begin{quote}
\textbf{That is the diagonal signal.} The bigram block asks a question the resampled table has
erased; the induction block asks one the sequence answers locally, and that answer survives the
average over $\pi$.
\end{quote}

\subsection{Reduction to a three-point correlation}

Expanding $m_u[w]=\frac1u\sum_{s<u}\1[v_s{=}w]$, the quantity to bound is
\begin{equation}
\E\bigl[\1[v_u{=}y_t]\,m_u[w]\bigr]
=\frac1u\sum_{s<u}\E\bigl[\1[v_u{=}y_t]\,\1[v_s{=}w]\bigr],\qquad s<u<t,
\end{equation}
so Statement~\ref{st:b} is a \emph{three-point correlation} in the chain, averaged over the random
table.

\subsection{Proposed proof structure}

Split the inner sum at a lag $L$.

\begin{description}
\item[Distant terms ($s\le u-L$).] Show
\begin{equation}
\bigl|\PP(v_s{=}w\mid v_u,v_t,y_t)-\PP(v_s{=}w)\bigr|\le C\rho^{\,u-s},
\end{equation}
so these contribute \emph{equally to every $w$} and cancel in the diagonal-minus-off-diagonal
contrast, leaving $O(\rho^L)$.
\item[Recent terms ($s>u-L$).] Dominated by $s=$ the last signal position before $u$. A one-step
Bayes computation gives a strictly positive constant on the diagonal, since occurrences of a
successor of $v$ are generated by transitions out of $v$.
\end{description}

\subsection{The two technical obstacles}

\begin{description}
\item[(i) Mixing over a random $2$-out graph.] Here $\rho$ is a mixing rate for the state chain.
For a \emph{fixed} table the chain need not mix well --- at $V_s=4$ one can draw tables with
transient states, so the stationary distribution is not uniform and $\rho$ is not uniformly
bounded. A \emph{high-probability statement over tables} is needed, together with control of the
exceptional set.
\item[(ii) The $1/u$ dilution.] $m_u$ averages the \emph{whole} prefix, so the per-position
diagonal boost is $O(1/u)$ while $m_u[w]$ itself is $O(1/V_s)$. The contrast survives the sums
over $t$ and $u$, but the constants should be explicit --- otherwise a referee will ask whether
the effect is real at the lengths actually trained on ($T=100$, scored on $[50,100)$).
\end{description}

%=====================================================================
\section{The readout bias \quad \PROVED}\label{sec:bias}

With a perfect distractor filter, $f_t$ is supported on $\Sig$ and sums to $1$, so
$\mathcal{Z}=(1+V_sb_s)+V_db_d$ is constant. Demanding that each distractor receive exactly
$(1-\alpha)/V_d$,
\begin{equation}
\frac{b_d}{\mathcal{Z}}=\frac{1-\alpha}{V_d}
\;\Longrightarrow\;V_db_d=(1-\alpha)\mathcal{Z}
\;\Longrightarrow\;\mathcal{Z}=(1+V_sb_s)+(1-\alpha)\mathcal{Z}
\;\Longrightarrow\;\mathcal{Z}=\frac{1+V_sb_s}{\alpha},
\end{equation}
whence
\begin{equation}
\boxed{\;b_d=\frac{(1-\alpha)(1+V_sb_s)}{V_d\,\alpha}\;}
\end{equation}
and the signal side receives $p_t[w]=\alpha\bigl(f_t[w]+b_s\bigr)/(1+V_sb_s)$, exactly
proportional to the smoothed count. The division by $\alpha$ is the part usually dropped: the
renormalisation shrinks whatever you put in.

\paragraph{Relation to a log-plus-softmax readout.}
If the nonlinearity before an unembedding is $z_w=\log(f_t[w]+b_s)$ on signal coordinates and a
constant $\beta$ on distractor coordinates, then
$\mathrm{softmax}(z)_w=(f_t[w]+b_s)/\sum_{w'}e^{z_{w'}}$, which is \eqref{eq:readout} with
$b_d=e^{\beta}$. The logarithm and the exponential cancel.
\NUM\ measured difference $5.6\times10^{-17}$.

\paragraph{Why the logarithm is not optional.}
Without it the softmax returns $e^{f_t[w]}$ --- exponential in the count rather than linear --- and
no rescaling repairs this, because $\mathrm{softmax}(\kappa f)$ is a different function of $f$
than $f/\sum f$ for every $\kappa$.

%=====================================================================
\section{Lemma 2: the first layer unlocks \quad \PROVED\ (modulo the Adam caveat)}

\begin{lemma}[After step $1$]\label{lem:step1}
Let $W^{(2)}_1=\eta G^{(2)}$ be the result of one gradient step from Lemma~\ref{lem:step0}, with
$G^{(2)}=-\partial L/\partial W^{(2)}$, and $W^{(1)}_1=0$ by Lemma~\ref{lem:step0}(i). Then:
\begin{enumerate}[label=(\roman*),nosep]
\item layer 1 is unchanged, so $h_t=m_t$ still;
\item with $\sigma_{t,u}=x^{(2)}_tG^{(2)}x^{(2)\top}_u$ and $\bar\sigma_t$ its mean over $u<t$,
\begin{equation}
a^{(2)}_{t,u}=\frac1t\Bigl(1+\frac{\eta}{\sqrt d}\bigl(\sigma_{t,u}-\bar\sigma_t\bigr)\Bigr)
+O(\eta^2);
\end{equation}
\item the first layer's gradient becomes non-zero, of order $\eta$:
\begin{equation}
-\frac{\partial L}{\partial W^{(1)}}\bigg|_{1}=\eta\,\Phi\bigl(G^{(2)}\bigr)+O(\eta^2),
\end{equation}
where $\Phi$ is linear in exactly those blocks of $G^{(2)}$ that contract with $h$, namely $Q$,
$R$ and $C$.
\end{enumerate}
\end{lemma}

\begin{proof}
(i) is Lemma~\ref{lem:step0}(i). (ii) is a first-order expansion of the softmax about the uniform
distribution. For (iii), the derivatives in the proof of Lemma~\ref{lem:step0}(i) are now evaluated
at $W^{(2)}=\eta G^{(2)}$ rather than $0$, so they equal $\eta$ times the corresponding blocks;
composing with $\partial h/\partial W^{(1)}$, which is $O(1)$ and non-zero, gives the claim. Only
the blocks multiplying $h$ survive the contraction.
\end{proof}

\begin{quote}
\textbf{Meaning.} The second layer \emph{recruits} the first. Layer 1 receives gradient only in
proportion to how much layer 2 has already learned to read $h$, and in the direction dictated by
the blocks that consume $h$ --- predominantly $G_Q$, whose diagonal structure is
Statement~\ref{st:b}.
\end{quote}

\paragraph{Caveat --- Adam.}
The experiments use Adam, whose first update has magnitude set by the learning rate almost
independently of the gradient scale. The $O(\eta)$ bookkeeping above is therefore the
gradient-flow / plain-GD statement. Either restate the scope, or repeat the measurement under
plain gradient descent.

%=====================================================================
\section{Lemma 3: what the first layer is pushed to compute \quad \TODO}

\begin{lemma}[After step $2$; intended]\label{lem:step2}
To leading order the update to $W^{(1)}$ increases $S^{(1)}_{u,s}$ for pairs $s<u$ with
$v_s\in\Sig$, decreases it for $v_s\in\Dis$, and among admissible $s$ the increase is monotone in
$s$. Consequently $h_u$ moves from the uniform prefix average $m_u$ towards the indicator of the
most recent signal symbol before $u$ --- that is, state recovery.
\end{lemma}

\noindent The intended argument:
\begin{enumerate}[nosep]
\item By Lemma~\ref{lem:step1}(iii) the layer-1 gradient is $\Phi(G^{(2)})$, dominated by $Q$.
\item $G_Q[v,\cdot]$ rewards keys $u$ whose $h_u$ places mass on $v$, when $v$ is the query symbol.
\item $h_u$ places mass on $v$ in proportion to how often $v$ occurs in $u$'s prefix; the update
therefore sharpens $h_u$ onto whichever symbol best predicts $v_u$, which by the generative process
is the state, i.e.\ the most recent signal symbol.
\item Distractor suppression and the recency ordering come from the two terms of the coefficient
(indicator minus baseline).
\end{enumerate}

\noindent This is the second piece of remaining work.

%=====================================================================
\section{Everything verified numerically}

All at $\alpha=0.6$ unless stated; single seed; exact full-batch gradients, not stochastic.

\begin{center}
\begin{tabular}{@{}ll>{\raggedright\arraybackslash}p{5.4cm}@{}}
\toprule
quantity & measured & relates to \\
\midrule
CE at zero initialisation & $2.039121$ & matches the unigram closed form \\
$\lVert\nabla_{W^{(1)}}L\rVert$ at zero init & $0.000\mathrm{e}{+}00$ exactly & Lemma~\ref{lem:step0}(i) \\
$G_P$ diag $-$ off-diag & $-1.475\times10^{-4}$ & Statement~\ref{st:a}, sign \\
$G_Q$ diag $-$ off-diag & $+9.621\times10^{-5}$ & Statement~\ref{st:b}, sign \\
$G_C$ diag $-$ off-diag & $+9.804\times10^{-6}$ & \\
$\cos(G_Q,\text{construction})$ & $+0.8670$ & random baseline $0.0099$ \\
$\cos(G_C,\text{construction})$ & $+0.8737$ & \\
self-successor rate & $49.95\%$ vs $2/V_s=50\%$ & Fact~\ref{fact:unif} \\
log+softmax $\equiv$ renormalisation & $5.6\times10^{-17}$ & Section~\ref{sec:bias} \\
post-distractor error, $\alpha\le0.7$ & $0.27\%$--$1.16\%$ & phase transition \\
post-distractor error, $\alpha=0.8$ & $51.59\%$ & phase transition \\
post-distractor error, $\alpha=0.9$ & $49.94\%$ & phase transition \\
\bottomrule
\end{tabular}
\end{center}

\noindent The phase-transition rows are measured on the $V=20$ models, $1500$ held-out sequences
each.

%=====================================================================
\section{What remains, in order}

\begin{enumerate}[nosep]
\item \textbf{Statement~\ref{st:a}} (Section~\ref{sec:a}). Self-contained, no mixing needed.
\emph{Start here.}
\item \textbf{Statement~\ref{st:b}} (Section~\ref{sec:b}). Needs the mixing bound over a random
$2$-out graph.
\item \textbf{Lemma~\ref{lem:step2}}.
\item \textbf{Scope.} The phase transition is measured on the $V=20$ models with a trained
unembedding and a log-plus-softmax readout; the dynamics above are for $V=8$ with the fixed
readout. These must be unified --- either replicate the transition at $V=8$, or re-derive the
dynamics for $V=20$.
\item \textbf{Seeds.} Every number above is a single seed.
\item \textbf{Post-signal control.} To claim that state recovery \emph{specifically} fails, the
post-signal error must stay flat while the post-distractor error explodes. This is not yet
measured correctly: the earlier attempt had an off-by-one in the state update, using the previous
signal symbol as the state at post-signal positions rather than $x_t$ itself.
\item \textbf{$\alpha\to1$ reduction.} Stated, neither proved nor measured.
\end{enumerate}

\end{document}
